\section{System Implementation}

\subsection{Frontend Implementation}

The frontend is implemented with Next.js and TypeScript. The app-router structure separates student and administrative workflows into route segments, while service modules in \code{frontend/src/services} centralize communication with the backend.

Important frontend concerns include:
\begin{itemize}
  \item authentication bootstrap and persistence;
  \item protected routes and permission checks;
  \item lesson and practice page composition;
  \item test flows with answer submission and result presentation;
  \item admin screens for users, roles, reports, and dashboards.
\end{itemize}

The component structure indicates a feature-oriented approach rather than a single monolithic page design. Examples include \code{QuestionPlayer.tsx}, \code{PracticeResultModal.tsx}, \code{AchievementGallery.tsx}, and \code{CreateTemplateModal.tsx}.

\subsection{Backend API Implementation}

The backend exposes versioned REST endpoints. Table \ref{tab:api-areas} groups the major API areas.

\begin{table}[H]
\centering
\caption{Main API groups}
\label{tab:api-areas}
\begin{tabularx}{\textwidth}{L{0.22\textwidth}Y}
\toprule
\textbf{API Group} & \textbf{Purpose} \\
\midrule
\code{/api/v1/auth} & Registration, login, profile retrieval, password change, and activity endpoints. \\
\code{/api/v1/lessons} & Lesson, grade, subject, mastery, practice initiation, and study-time endpoints. \\
\code{/api/v1/tests} & Grade tests, attempt detail, answer submission, finish flow, template management, and report workflow. \\
\code{/api/v1/admin} & Users, roles, access control, insights, and dashboard statistics. \\
\code{/api/v1/achievements} & Achievement seeding, listing, earned-state retrieval, and re-checking. \\
\bottomrule
\end{tabularx}
\end{table}

The API design is pragmatic. It is expressive enough for the frontend and explicit enough for debugging, while remaining simple to secure and deploy.

\subsection{Question Generation and Grading}

The file \code{backend/services/mathService.ts} contains a central part of the project: logic-driven variable generation and answer evaluation. The implementation uses the \code{mathjs} expression engine for parsing and evaluating formulas, derived expressions, and constraints \cite{mathjs-docs}. The service supports several rule types:
\begin{itemize}
  \item fixed numeric values;
  \item arrays of possible values;
  \item numeric ranges with optional step and precision;
  \item random choices from configured lists;
  \item derived expressions evaluated from previously generated variables;
  \item constraints that must hold before a question is accepted.
\end{itemize}

The generation flow can be summarized as:
\begin{enumerate}
  \item parse the stored \code{logic\_config};
  \item generate variables from configured rules;
  \item apply derived expressions;
  \item verify constraints;
  \item repeat until a valid assignment is found or a maximum number of attempts is reached.
\end{enumerate}

Answer checking supports multiple cases:
\begin{itemize}
  \item numeric answers, compared with a small tolerance;
  \item boolean answers, normalized from common textual forms;
  \item text or structured answers, normalized for comparison;
  \item multiple accepted formulas for the same question.
\end{itemize}

This is a stronger implementation than a naive string-equality checker because it tolerates common input variance while still grounding correctness in mathematically evaluated formulas.

\subsection{Attempt Lifecycle}

The implemented attempt lifecycle is one of the most important interactions in the system.

\subsubsection*{Practice Start}

When a student starts practice for a lesson:
\begin{enumerate}
  \item the backend validates lesson existence and subject access;
  \item it loads lesson templates;
  \item it generates variables and right answers for selected templates;
  \item it creates a \code{test\_attempt} record;
  \item it bulk-inserts \code{question\_snapshots};
  \item it returns a stable attempt payload to the frontend.
\end{enumerate}

\subsubsection*{Answer Submission}

When a student submits an answer:
\begin{enumerate}
  \item the backend loads the snapshot and template;
  \item it evaluates the answer using either text comparison or formula-based checking;
  \item it updates the stored snapshot with answer and correctness state;
  \item it returns correctness information and explanation content.
\end{enumerate}

\subsubsection*{Attempt Completion}

When an attempt is finished:
\begin{enumerate}
  \item unanswered snapshots are marked incorrect;
  \item correct answers are counted and the final score is computed;
  \item lesson mastery is updated when the attempt is linked to a lesson;
  \item XP is awarded and logged;
  \item aggregate student statistics are recalculated;
  \item achievements are re-evaluated and may award additional XP.
\end{enumerate}

This completion pipeline shows that the project treats scoring as part of a broader progression workflow rather than as an isolated calculation.

\subsection{Administrative Implementation}

The administrative subsystem includes user and role management, subject permission management, dashboard statistics, and question report review. The route design suggests that administrators can change both coarse and fine-grained access behavior without redeploying the application.

This is reinforced by the underlying data model:
\begin{itemize}
  \item \code{roles} store named roles;
  \item \code{actions} and \code{resources} define the authorization vocabulary;
  \item \code{role\_permissions} map actions to resources per role;
  \item \code{role\_subject\_permissions} restrict which subjects a role can manage.
\end{itemize}

This structure scales better than a hard-coded permission matrix embedded directly in frontend components.

\subsection{Configuration and Deployment}

The repository defines environment variables for both frontend and backend execution. Examples include:
\begin{itemize}
  \item API base URLs and rewrite targets for the frontend;
  \item \code{DATABASE\_URL}, \code{JWT\_SECRET}, and \code{CORS\_ORIGINS} for the backend;
  \item \code{AUTO\_SEED\_ACHIEVEMENTS} to control whether achievement definitions are seeded automatically.
\end{itemize}

The backend startup process explicitly checks the achievement seeding flag before inserting seed data. This is a sound operational choice because it prevents accidental reseeding in every environment.
